\part{SYSTEMATIC DEDUCTION OF FUNDAMENTAL LAWS OF NATURE}\label{part:2}

\section{Emergent Gravity, Quantum Nature of the Equivalence Principle, and Horizon Thermodynamics}\label{sec:rozdil2}

\subsection{Massless Spin-2 Gauge Field and Non-Linear Gravitational Self-Interaction (Feynman-Deser Theorem)}\label{sec:2.1}

\subsubsection{Massless Spin-2 Gauge Boson in Minkowski Spacetime}\label{sec:2.1.1}

In classical relativistic field theory, the gravitational field is traditionally postulated via a fundamental geometric curvature of the spacetime continuum. In our contour approach, pseudo-Riemannian geometry and gravitation are \textbf{emergent macroscopic consequences} of the interaction of a massless rank-2 gauge field in a flat vacuum.

Consider a massless symmetric rank-2 tensor field $h_{\mu\nu} = h_{\nu\mu}$ in Minkowski spacetime with the flat metric tensor $\eta_{\mu\nu} = \text{diag}(+1, -1, -1, -1)$. The quadratic free Lagrangian of Markus Fierz and Wolfgang Pauli (1939) possesses the unique structure:

\begin{equation}\label{eq:2.1}\tag{2.1}
\mathcal{L}_{\text{FP}} = \frac{1}{2} \partial_\lambda h_{\mu\nu} \partial^\lambda h^{\mu\nu} - \partial_\mu h^{\mu\nu} \partial_\lambda h^\lambda_{\;\nu} + \partial_\mu h^{\mu\nu} \partial_\nu h - \frac{1}{2} \partial_\mu h \partial^\mu h,
\end{equation}

where $h_{\mu\nu}$ is the dynamical symmetric rank-2 tensor field (graviton);  
$h \equiv \eta^{\mu\nu} h_{\mu\nu}$ is the Lorentz trace of the field tensor;  
$\partial_\mu = \frac{\partial}{\partial x^\mu}$ is the partial 4-derivative operator.

The Fierz-Pauli Lagrangian (2.1) is strictly invariant under Abelian gradient gauge transformations:

\begin{equation}\label{eq:2.2}\tag{2.2}
h_{\mu\nu}(x) \to h'_{\mu\nu}(x) = h_{\mu\nu}(x) + \partial_\mu \xi_\nu(x) + \partial_\nu \xi_\mu(x),
\end{equation}

where $\xi_\mu(x)$ is an arbitrary differentiable vector gauge field of local coordinate shifts.

The free field equations of motion are obtained by varying the action $S_{\text{FP}} = \int d^4x \mathcal{L}_{\text{FP}}$ with respect to the field $h^{\mu\nu}$:

\begin{equation}\label{eq:2.3}\tag{2.3}
\mathcal{E}_{\mu\nu}^{\alpha\beta} h_{\alpha\beta} = 0,
\end{equation}

where $\mathcal{E}_{\mu\nu}^{\alpha\beta}$ is the self-adjoint differential tensor wave operator of Fierz and Pauli:

\begin{equation}\label{eq:2.4}\tag{2.4}
\mathcal{E}_{\mu\nu}^{\alpha\beta} h_{\alpha\beta} = \square h_{\mu\nu} - \partial_\mu \partial_\alpha h^\alpha_{\;\nu} - \partial_\nu \partial_\alpha h^\alpha_{\;\mu} + \partial_\mu \partial_\nu h + \eta_{\mu\nu} \left( \partial_\alpha \partial_\beta h^{\alpha\beta} - \square h \right).
\end{equation}

Introducing the trace-reversed gauge field $\bar{h}_{\mu\nu} \equiv h_{\mu\nu} - \frac{1}{2}\eta_{\mu\nu} h$ and imposing the harmonic Lorenz-Hilbert gauge condition:

\begin{equation}\label{eq:2.5}\tag{2.5}
\partial^\mu \bar{h}_{\mu\nu} = 0,
\end{equation}

equation (2.3) simplifies to the linear d'Alembert wave equation:

\begin{equation}\label{eq:2.6}\tag{2.6}
\square \bar{h}_{\mu\nu} = 0.
\end{equation}

Due to the gauge invariance (2.2), out of the 10 independent components of the symmetric tensor $h_{\mu\nu}$, exactly 4 components are eliminated by the gauge choice (2.5), while residual gauge transformations with $\square \xi_\mu = 0$ remove another 4 unphysical degrees of freedom. The physical field of the massless graviton possesses precisely $10 - 4 - 4 = \mathbf{2}$ dynamical degrees of freedom of transverse polarization with spin 2 (the $h_+$ and $h_\times$ polarizations).

\subsubsection{Non-Linear Gravitational Self-Interaction: Derivation of the Einstein-Hilbert Action}\label{sec:2.1.2}

As was first realized by Robert Kraichnan \cite{ref13} (1955), Richard Feynman (1962), Stanley Deser \cite{ref14} (1970), and Steven Weinberg (1965), the linear Fierz-Pauli theory is inherently incomplete and inconsistent. The reason for this is that the quantum gravitational field $h_{\mu\nu}$ carries energy, momentum, and stress, and therefore, by the universal law of nature, must gravitationally interact \textbf{with itself}.

Let the field $h_{\mu\nu}$ at the first step couple linearly to the symmetric energy-momentum tensor of matter $T_{\mu\nu}^{(0)}$:

\begin{equation}\label{eq:2.7}\tag{2.7}
\mathcal{L}^{(1)} = \mathcal{L}_{\text{FP}} + \frac{1}{2}\kappa \, h_{\mu\nu} T^{\mu\nu}_{(0)},
\end{equation}

where $\kappa = \sqrt{32\pi G / c^4}$ is the gravitational coupling constant.

However, the field equations obtained by varying the Lagrangian (2.7) take the form:

\begin{equation}\label{eq:2.8}\tag{2.8}
\mathcal{E}_{\mu\nu}^{\alpha\beta} h_{\alpha\beta} = -\frac{1}{2}\kappa T_{\mu\nu}^{(0)}.
\end{equation}

Taking the 4-divergence $\partial^\mu$ of both sides of equality (2.8), the left-hand side vanishes identically by virtue of the Noether identity stemming from gauge invariance ($\partial^\mu \mathcal{E}_{\mu\nu}^{\alpha\beta} h_{\alpha\beta} \equiv 0$). However, the divergence of the matter tensor in the presence of a gravitational field is \textbf{non-zero}:

\begin{equation}\label{eq:2.9}\tag{2.9}
\partial^\mu T_{\mu\nu}^{(0)} \ne 0.
\end{equation}

A contradiction arises: the left-hand side of the equation is conserved, while the right-hand side is not. To restore current conservation, one must add the field's own canonical energy-momentum tensor of gravitons, $t_{\mu\nu}^{(1)}(h)$, to the source:

\begin{equation}\label{eq:2.10}\tag{2.10}
\mathcal{E}_{\mu\nu}^{\alpha\beta} h_{\alpha\beta} = -\frac{1}{2}\kappa \left( T_{\mu\nu}^{(0)} + t_{\mu\nu}^{(1)}(h) \right).
\end{equation}

This requires adding a cubic field self-interaction term $\kappa \mathcal{L}^{(3)}(h^3)$ to the original Lagrangian:

\begin{equation}\label{eq:2.11}\tag{2.11}
\mathcal{L}^{(2)} = \mathcal{L}_{\text{FP}} + \kappa \mathcal{L}^{(3)}(h^3) + \frac{1}{2}\kappa \, h_{\mu\nu} T^{\mu\nu}.
\end{equation}

However, this new cubic term generates a new second-order contribution in $\kappa$ to the energy-momentum tensor, requiring the introduction of a quartic term $\kappa^2 \mathcal{L}^{(4)}(h^4)$, and so on ad infinitum:

\begin{equation}\label{eq:2.12}\tag{2.12}
\mathcal{L}_{\text{total}} = \mathcal{L}_{\text{FP}} + \sum_{n=1}^\infty \kappa^n \mathcal{L}^{(n+2)}(h^{n+2}).
\end{equation}

\subsubsection{Feynman-Deser Theorem}\label{sec:feynman-deser-thm}

\begin{theorem}[Feynman-Deser Theorem]\label{thm:feynman-deser}
The infinite non-linear self-interaction series of a massless rank-2 gauge field (2.12) sums rigorously and uniquely into the closed geometric Einstein-Hilbert action of a pseudo-Riemannian manifold:

\begin{equation}\label{eq:2.13}\tag{2.13}
S_{\text{total}}[g] = \frac{c^4}{16\pi G} \int d^4x \sqrt{-g} \, R,
\end{equation}

where the full non-linear spacetime metric $g_{\mu\nu}(x)$ is determined as a non-linear conformal matrix exponential of perturbations:

\begin{equation}\label{eq:2.14}\tag{2.14}
g_{\mu\nu}(x) = \eta_{\mu\nu} + \kappa h_{\mu\nu}(x) + \frac{1}{4}\kappa^2 h_{\mu\alpha}(x) h^\alpha_{\;\nu}(x) + \dots.
\end{equation}
\end{theorem}

The Feynman-Deser theorem is a landmark result of theoretical physics: \textbf{spacetime curvature and Einstein's General Relativity are not an arbitrary geometry, but an inevitable consequence of the quantum self-interaction of spin-2 contour phase excitations}.

\subsection{Variational Derivation of Einstein Field Equations and Reduction to Classical GR in the Strong-Field Regime}\label{sec:2.2}

Consider the variation of the total action of the system, consisting of the Einstein-Hilbert gravitational action $S_{\text{EH}}$ and the matter field action $S_{\text{matter}}$:

\begin{equation}\label{eq:2.15}\tag{2.15}
S_{\text{total}}[g, \psi] = S_{\text{EH}}[g] + S_{\text{matter}}[g, \psi] = \frac{c^4}{16\pi G} \int d^4x \sqrt{-g} \, R + \int d^4x \sqrt{-g} \, \mathcal{L}_{\text{matter}}(g, \psi).
\end{equation}

Varying the total action with respect to the contravariant metric tensor $g^{\mu\nu}(x)$:

\begin{equation}\label{eq:2.16}\tag{2.16}
\delta S_{\text{total}} = \int d^4x \left( \frac{\delta S_{\text{EH}}}{\delta g^{\mu\nu}} + \frac{\delta S_{\text{matter}}}{\delta g^{\mu\nu}} \right) \delta g^{\mu\nu} = 0.
\end{equation}

We compute the variation of the Einstein-Hilbert action density $\delta(\sqrt{-g} R)$:

\begin{equation}\label{eq:2.17}\tag{2.17}
\delta(\sqrt{-g} R) = (\delta\sqrt{-g}) R + \sqrt{-g} (\delta R).
\end{equation}

Using Jacobi's formula for the variation of the metric determinant:

\begin{equation}\label{eq:2.18}\tag{2.18}
\delta g = g \, g^{\mu\nu} \delta g_{\mu\nu} = -g \, g_{\mu\nu} \delta g^{\mu\nu} \implies \delta\sqrt{-g} = -\frac{1}{2}\sqrt{-g} \, g_{\mu\nu} \delta g^{\mu\nu}.
\end{equation}

The variation of the Ricci scalar $R = g^{\mu\nu} R_{\mu\nu}$ contains two terms:

\begin{equation}\label{eq:2.19}\tag{2.19}
\delta R = R_{\mu\nu} \delta g^{\mu\nu} + g^{\mu\nu} \delta R_{\mu\nu}.
\end{equation}

By Palatini's identity, the variation of the Ricci tensor is expressed through the covariant derivative of the variation of the Christoffel connection $\delta\Gamma^\lambda_{\mu\nu}$:

\begin{equation}\label{eq:2.20}\tag{2.20}
\delta R_{\mu\nu} = \nabla_\lambda (\delta\Gamma^\lambda_{\mu\nu}) - \nabla_\nu (\delta\Gamma^\lambda_{\mu\lambda}).
\end{equation}

Then the second term in (2.19) contracts into a total 4-divergence of a vector:

\begin{equation}\label{eq:2.21}\tag{2.21}
\sqrt{-g} \, g^{\mu\nu} \delta R_{\mu\nu} = \sqrt{-g} \, \nabla_\mu \left( g^{\alpha\beta} \delta\Gamma^\mu_{\alpha\beta} - g^{\mu\alpha} \delta\Gamma^\beta_{\alpha\beta} \right) = \partial_\mu \left( \sqrt{-g} \, W^\mu \right).
\end{equation}

By the Gauss-Ostrogradsky theorem, the integral of a total divergence transforms into a surface integral at the asymptotic boundary of spacetime and vanishes for fixed metric variations on the boundary.

Substituting results (2.18) and (2.19) into the variation of the gravitational action:

\begin{equation}\label{eq:2.22}\tag{2.22}
\frac{\delta S_{\text{EH}}}{\delta g^{\mu\nu}} = \frac{c^4}{16\pi G} \sqrt{-g} \left( R_{\mu\nu} - \frac{1}{2} g_{\mu\nu} R \right).
\end{equation}

Defining the metric symmetric energy-momentum tensor of matter according to David Hilbert:

\begin{equation}\label{eq:2.23}\tag{2.23}
T_{\mu\nu} \equiv -\frac{2}{\sqrt{-g}} \frac{\delta S_{\text{matter}}}{\delta g^{\mu\nu}} = -2 \frac{\partial \mathcal{L}_{\text{matter}}}{\partial g^{\mu\nu}} + g_{\mu\nu} \mathcal{L}_{\text{matter}},
\end{equation}

where $T_{\mu\nu}$ is the energy-momentum tensor of matter fields.

Setting the total variation (2.16) to zero for arbitrary variations $\delta g^{\mu\nu}$, we obtain the classical relativistic \textbf{Einstein field equations}:

\begin{equation}\label{eq:2.24}\tag{2.24}
G_{\mu\nu} \equiv R_{\mu\nu} - \frac{1}{2} g_{\mu\nu} R = \frac{8\pi G}{c^4} T_{\mu\nu},
\end{equation}

where $G_{\mu\nu}$ is the symmetric rank-2 Einstein tensor.

Due to the geometric differential Bianchi identities $\nabla^\mu G_{\mu\nu} \equiv 0$, the covariant divergence of the left-hand side of equations (2.24) vanishes identically. This automatically guarantees the fulfillment of the fundamental physical \textbf{covariant conservation law of energy and momentum of matter} in curved spacetime:

\begin{equation}\label{eq:2.25}\tag{2.25}
\nabla^\mu T_{\mu\nu} = 0.
\end{equation}

Under the dominant energy condition for the tensor $T_{\mu\nu}$, an asymptotically flat gravitational configuration is characterized by a strictly positive total mass (ADM energy), as proven in the positive energy theorem by Richard Schoen and Shing-Tung Yau \cite{ref15}.

\subsubsection{Reduction to Classical GR in the Strong Field and Automatic Satisfaction of Post-Newtonian Tests}\label{sec:2.2.1}

Since in the regime of strong and moderate accelerations ($a \gg a_0$, where $a_0 \approx 1.2 \cdot 10^{-10}\text{ m}\cdot\text{s}^{-2}$ is the quantum dephasing threshold) the vacuum interpolation function tends to unity:

\begin{equation*}
\lim_{a/a_0 \to \infty} \mu\left(\frac{a}{a_0}\right) = 1,
\end{equation*}

the dynamical field equations of the phase contour theory identically coincide with the classical non-linear tensor Einstein equations of GR:

\begin{equation*}
G_{\mu\nu} = \frac{8\pi G}{c^4} T_{\mu\nu}.
\end{equation*}

As a consequence of this exact identity, all classical vacuum solutions of GR (the spherically symmetric Schwarzschild geometry, the rotating Kerr metric, the electrostatic Reissner-Nordstr\"om solution) are rigorous solutions of the developed theory at the corresponding scale.

Accordingly, all textbook experimental and post-Newtonian tests of GR:
\begin{enumerate}
\item \textbf{Relativistic perihelion precession of Mercury} ($\Delta\phi = \frac{6\pi GM}{a(1-e^2)c^2} \approx 42.98''/\text{century}$);
\item \textbf{Gravitational deflection of light by massive bodies} ($\theta = \frac{4GM}{R c^2} \approx 1.751''$ for a ray grazing the solar limb);
\item \textbf{Shapiro gravitational time delay effect} ($\Delta t = \frac{4GM}{c^3} \left[ \ln\left(\frac{4 r_E r_R}{R_0^2}\right) + 1 \right]$);
\item \textbf{Quadrupole gravitational wave radiation and binary pulsars} ($P_{\text{GW}} = -\frac{G}{5 c^5} \langle \dddot{I}_{ij} \dddot{I}^{ij} \rangle$, confirmed by LIGO/Virgo observations);
\end{enumerate}
are satisfied automatically with the same limiting precision as in standard General Relativity. This eliminates the need for their separate calculation and focuses our investigation on regimes where standard GR faces a crisis (ultra-weak accelerations, dark matter, and quantum decoherence).

\subsection{Non-Abelian Yang-Mills Gauge Dynamics, Wilson Loops, and Contour Origin of Hadron Mass}\label{sec:2.3}

\subsubsection{Reduction of the Universal Master-Functional onto the Non-Abelian Subgroup $SU(3)_c$}\label{sec:2.3.1}

In Section~\hyperref[sec:1.2]{1.2}, the fundamental state of the quantum Universe was defined via the universal Richard Feynman contour functional on a spacelike 3-hypersurface $\Sigma$:

\begin{equation}\label{eq:2.3.1}\tag{2.3.1}
\mathbf{\Psi}[\Sigma] = \int_{\mathcal{F}_\Sigma} \mathcal{D}[\gamma] \, \Delta_{\text{FP}}[\gamma] \, \frac{1}{\mathcal{N}} \text{Tr}\left[ \mathcal{P}\exp\left( \frac{i}{\hbar} \oint_{\gamma} \mathbf{\Pi}_\mu dx^\mu \right) \right],
\end{equation}

where the unified 1-form operator of the generalized canonical 4-momentum $\mathbf{\Pi}_\mu$ acts in the total space of the principal bundle $\mathcal{P}(\mathcal{M}, \mathcal{G})$ with structural Lie algebra $\mathfrak{g} = \mathfrak{u}(1) \oplus \mathfrak{su}(2) \oplus \mathfrak{su}(3) \oplus \mathfrak{so}(3,1)$ of dimension $\mathcal{N} = 1 + 3 + 8 + 6 = 18$ (with an exact fundamental matrix representation of dimension $\mathcal{N} = 10$).

We perform an exact algebraic projection of the operator $\mathbf{\Pi}_\mu$ onto the non-Abelian chromodynamic component $\mathfrak{su}(3)_c$, responsible for the strong nuclear interaction. The corresponding sector of the momentum 1-form factorizes as:

\begin{equation}\label{eq:2.3.2}\tag{2.3.2}
\mathbf{\Pi}_\mu \Big|_{\mathfrak{su}(3)} = p_\mu \, \mathbf{I}_3 + g \, A_\mu^a T^a,
\end{equation}

where $p_\mu$ is the classical 4-momentum of translational motion along the contour $\gamma$;  
$\mathbf{I}_3$ is the identity operator in the fundamental three-dimensional color space $\mathbb{C}^3$;  
$g$ is the dimensionless fundamental Yang-Mills gauge coupling constant;  
$A_\mu^a(x)$ are the non-Abelian gauge potentials (gluon field vector potentials);  
$T^a$ ($a = 1, \dots, 8$) are the Hermitian traceless generators of the fundamental representation of the Lie algebra $\mathfrak{su}(3)$, normalized by the standard relation:

\begin{equation}\label{eq:2.3.3}\tag{2.3.3}
\text{Tr}\left( T^a T^b \right) = \frac{1}{2} \delta^{ab},
\end{equation}

and satisfying the non-Abelian commutation algebra:

\begin{equation}\label{eq:2.3.4}\tag{2.3.4}
\left[ T^a, T^b \right] = i f^{abc} T^c,
\end{equation}

where $f^{abc}$ are the completely antisymmetric structure constants of the color group $SU(3)$.

Substituting the projection (2.3.2) into the closed contour phase accumulation of the functional (2.3.1), we find that along a closed spacetime loop $\gamma: [0, 1] \to \mathcal{M}$, the phase factor decomposes into an Abelian kinematic phase and a matrix Wilson loop operator:

\begin{equation}\label{eq:2.3.5}\tag{2.3.5}
\mathbf{W}[\gamma] = \mathcal{P}\exp\left( \frac{ig}{\hbar} \oint_{\gamma} A_\mu^a(x) T^a dx^\mu \right) \in SU(3)_c.
\end{equation}

Here the Dyson path-ordering operator $\mathcal{P}$ acts along the monotonic contour parameter $s \in [0, 1]$ ($x^\mu = x^\mu(s)$, $x^\mu(0) = x^\mu(1)$):

\begin{equation}\label{eq:2.3.6}\tag{2.3.6}
\mathcal{P}\left[ A(s_1) A(s_2) \right] = \theta(s_1 - s_2) A(s_1) A(s_2) + \theta(s_2 - s_1) A(s_2) A(s_1).
\end{equation}

\subsubsection{Non-Abelian Stokes Theorem and Derivation of the Field Strength Tensor $F_{\mu\nu}^a$}\label{sec:2.3.2}

Consider a closed contour $\delta\gamma$ encircling an infinitesimal oriented 2-plaquette $d\sigma^{\mu\nu}$ centered at the point $x_0^\mu$. We parameterize the contour as a rectangular parallelogram with vertices $x_0$, $x_0 + \Delta x^\mu$, $x_0 + \Delta x^\mu + \Delta x^\nu$, and $x_0 + \Delta x^\nu$.

We expand the Wilson loop operator $\mathbf{W}[\delta\gamma]$ in a Taylor series in powers of the oriented area bivector $d\sigma^{\mu\nu} = \Delta x^\mu \wedge \Delta x^\nu$:

\begin{equation}\label{eq:2.3.7}\tag{2.3.7}
\mathbf{W}[\delta\gamma] = \mathbf{I}_3 + \frac{ig}{\hbar} \oint_{\delta\gamma} A_\mu(x) dx^\mu - \frac{g^2}{2\hbar^2} \mathcal{P} \left( \oint_{\delta\gamma} A_\mu(x) dx^\mu \right)^2 + \mathcal{O}(d\sigma^{3/2}).
\end{equation}

Applying Taylor's formula for the gauge connection $A_\mu(x) = A_\mu^a(x) T^a$:

\begin{equation*}
A_\mu(x_0 + \Delta x) = A_\mu(x_0) + \Delta x^\alpha \partial_\alpha A_\mu(x_0) + \frac{1}{2} \Delta x^\alpha \Delta x^\beta \partial_\alpha \partial_\beta A_\mu(x_0) + \dots
\end{equation*}

Computing the integration along four consecutive segments of the closed circuit yields:
\begin{enumerate}
\item Segment 1 ($x_0 \to x_0 + \Delta x^\mu$): $A_\mu(x_0 + \frac{1}{2}\Delta x^\mu) \Delta x^\mu$;
\item Segment 2 ($x_0 + \Delta x^\mu \to x_0 + \Delta x^\mu + \Delta x^\nu$): $A_\nu(x_0 + \Delta x^\mu + \frac{1}{2}\Delta x^\nu) \Delta x^\nu$;
\item Segment 3 ($x_0 + \Delta x^\mu + \Delta x^\nu \to x_0 + \Delta x^\nu$): $- A_\mu(x_0 + \frac{1}{2}\Delta x^\mu + \Delta x^\nu) \Delta x^\mu$;
\item Segment 4 ($x_0 + \Delta x^\nu \to x_0$): $- A_\nu(x_0 + \frac{1}{2}\Delta x^\nu) \Delta x^\nu$.
\end{enumerate}

The sum of first derivatives at linear order in the area $\Delta x^\mu \Delta x^\nu$ determines the curl of the Abelian part:

\begin{equation}\label{eq:2.3.8}\tag{2.3.8}
\oint_{\delta\gamma} A_\alpha dx^\alpha = \left( \partial_\mu A_\nu - \partial_\nu A_\mu \right) \Delta x^\mu \Delta x^\nu.
\end{equation}

The second term of the Dyson ordered path integral expansion contains the quadratic commutator of operators at different points of the loop:

\begin{equation}\label{eq:2.3.9}\tag{2.3.9}
- \frac{g^2}{2\hbar^2} \mathcal{P} \left( \oint A_\mu dx^\mu \right)^2 \implies - \frac{ig}{\hbar} \left( - \frac{ig}{\hbar} [A_\mu, A_\nu] \right) \Delta x^\mu \Delta x^\nu = - \frac{g^2}{\hbar^2} [A_\mu, A_\nu] \Delta x^\mu \Delta x^\nu.
\end{equation}

Combining both contributions, we obtain the exact non-Abelian curvature of the phase contour:

\begin{equation}\label{eq:2.3.10}\tag{2.3.10}
\mathbf{W}[\delta\gamma] = \mathbf{I}_3 + \frac{ig}{2\hbar} \mathbf{F}_{\mu\nu}(x_0) \, d\sigma^{\mu\nu} - \frac{g^2}{8\hbar^2} \mathbf{F}_{\mu\nu}(x_0) \mathbf{F}_{\alpha\beta}(x_0) \, d\sigma^{\mu\nu} d\sigma^{\alpha\beta} + \dots,
\end{equation}

where the Yang-Mills field strength 2-form $\mathbf{F}_{\mu\nu} = F_{\mu\nu}^a T^a$ is defined as the commutator of gauge-covariant derivatives $D_\mu = \partial_\mu - \frac{ig}{\hbar} A_\mu^a T^a$:

\begin{equation}\label{eq:2.3.11}\tag{2.3.11}
\mathbf{F}_{\mu\nu} \equiv \frac{\hbar}{ig} \left[ D_\mu, D_\nu \right] = \partial_\mu \mathbf{A}_\nu - \partial_\nu \mathbf{A}_\mu - \frac{ig}{\hbar} \left[ \mathbf{A}_\mu, \mathbf{A}_\nu \right].
\end{equation}

Projecting onto the generators $T^a$ using the commutation relation (2.3.4), we obtain the classical form of the components of the Yang-Mills field strength tensor:

\begin{equation}\label{eq:2.3.12}\tag{2.3.12}
F_{\mu\nu}^a = \partial_\mu A_\nu^a - \partial_\nu A_\mu^a + \frac{g}{\hbar} f^{abc} A_\mu^b A_\nu^c.
\end{equation}

\subsubsection{Variational Derivation of the Non-Abelian Yang-Mills Action}\label{sec:2.3.3}

We take the gauge trace ($\text{Tr}$) of the contour operator (2.3.10) in the fundamental representation. Taking into account the tracelessness of the generators $\text{Tr}(T^a) = 0$, the linear term in the area vanishes identically:

\begin{equation}\label{eq:2.3.13}\tag{2.3.13}
\text{Tr}\left( \mathbf{F}_{\mu\nu} \right) = F_{\mu\nu}^a \, \text{Tr}(T^a) \equiv 0.
\end{equation}

The leading non-vanishing contribution to the contour functional arises at second order in the expansion with respect to the area of the closed loop:

\begin{equation}\label{eq:2.3.14}\tag{2.3.14}
\text{Tr}\left( \mathbf{W}[\delta\gamma] \right) = \text{Tr}(\mathbf{I}_3) - \frac{g^2}{8\hbar^2} \text{Tr}\left( \mathbf{F}_{\mu\nu} \mathbf{F}_{\alpha\beta} \right) d\sigma^{\mu\nu} d\sigma^{\alpha\beta} + \dots
\end{equation}

Since $\text{Tr}(\mathbf{I}_3) = 3$, and according to the normalization (2.3.3):

\begin{equation}\label{eq:2.3.15}\tag{2.3.15}
\text{Tr}\left( \mathbf{F}_{\mu\nu} \mathbf{F}_{\alpha\beta} \right) = F_{\mu\nu}^a F_{\alpha\beta}^b \, \text{Tr}(T^a T^b) = \frac{1}{2} F_{\mu\nu}^a F_{\alpha\beta}^a.
\end{equation}

For isotropic statistical averaging over all possible spacetime orientations of random microscopic Feynman loops in a Euclidean or pseudo-Riemannian volume $d^4x$:

\begin{equation}\label{eq:2.3.16}\tag{2.3.16}
\langle d\sigma^{\mu\nu} d\sigma^{\alpha\beta} \rangle = \frac{1}{3} \left( g^{\mu\alpha} g^{\nu\beta} - g^{\mu\beta} g^{\nu\alpha} \right) (d\sigma)^2.
\end{equation}

Contracting the indices, we obtain the scalar:

\begin{equation}\label{eq:2.3.17}\tag{2.3.17}
F_{\mu\nu}^a F_{\alpha\beta}^a \left( g^{\mu\alpha} g^{\nu\beta} - g^{\mu\beta} g^{\nu\alpha} \right) = 2 F_{\mu\nu}^a F^{\mu\nu}_a.
\end{equation}

Consequently, the exponential accumulation of the effective contour phase of an ensemble of closed non-Abelian Feynman loops generates the action integral:

\begin{equation}\label{eq:2.3.18}\tag{2.3.18}
\mathbf{\Psi}_{\text{YM}} \sim \exp\left( \frac{i}{\hbar} S_{\text{YM}}[A] \right),
\end{equation}

where $S_{\text{YM}}$ is the fundamental classical Yang-Mills action:

\begin{equation}\label{eq:2.3.19}\tag{2.3.19}
S_{\text{YM}}[A] = -\frac{1}{4} \int d^4x \, F_{\mu\nu}^a F^{\mu\nu}_a.
\end{equation}

Varying the action with respect to the gauge potentials, $\frac{\delta S_{\text{YM}}}{\delta A_\nu^a} = 0$, yields the non-linear dynamical Yang-Mills field equations:

\begin{equation}\label{eq:2.3.20}\tag{2.3.20}
D_\mu F^{\mu\nu}_a \equiv \partial_\mu F^{\mu\nu}_a + \frac{g}{\hbar} f^{abc} A_\mu^b F^{\mu\nu}_c = 0.
\end{equation}

\begin{theorem}[Theorem 2.3.1 (Contour Emergence of Yang-Mills)]\label{thm:2.3.1}
Non-Abelian Yang-Mills gauge fields are not a priori postulated fundamental entities. They are the direct phase curvature of Feynman contour trajectories in the corresponding gauge bundle $\mathcal{P}(\mathcal{M}, SU(3))$, and the classical Yang-Mills equations (2.3.20) emerge as the condition of stationary phase of the Master-Functional $\mathbf{\Psi}[\Sigma]$.
\end{theorem}

\subsubsection{Non-Linear Loop Corrections, Phase Vacuum Polarization, and Asymptotic Freedom}\label{sec:2.3.4}

Consider quantum fluctuations of contour trajectories around the classical configuration vacuum: $A_\mu^a \to A_\mu^a + a_\mu^a$, where $a_\mu^a$ is the quantum perturbation field over which the path integration in the functional $\mathbf{\Psi}[\Sigma]$ is performed.

Due to the non-Abelian nature of the commutator $[T^a, T^b] \neq 0$, the action Lagrangian (2.3.19) contains three-point and four-point self-interaction vertices of gluon trajectories:

\begin{equation}\label{eq:2.3.21}\tag{2.3.21}
\mathcal{L}_{\text{int}} = - \frac{g}{\hbar} f^{abc} (\partial_\mu A_\nu^a) A^{\mu b} A^{\nu c} - \frac{g^2}{4\hbar^2} f^{abc} f^{ade} A_\mu^b A_\nu^c A^{\mu d} A^{\nu e}.
\end{equation}

The computation of the one-loop vacuum polarization operator $\Pi_{\mu\nu}^{ab}(k)$ for a non-Abelian gauge field in the Feynman gauge includes three types of closed loops:
\begin{enumerate}
\item A gluon loop with two triple vertices;
\item A gluon contact tadpole/seagull loop with a quartic vertex;
\item A loop of Faddeev-Popov ghosts $c^a, \bar{c}^a$, originating from the determinant $\Delta_{\text{FP}}[\gamma]$ and ensuring the cancellation of unphysical timelike and longitudinal gauge path modes.
\end{enumerate}

\begin{equation}\label{eq:2.3.22}\tag{2.3.22}
\Pi_{\mu\nu}^{ab}(k) = \delta^{ab} \left( k^2 g_{\mu\nu} - k_\mu k_\nu \right) \Pi(k^2).
\end{equation}

Regularization of divergences via the 't~Hooft-Veltman dimensional regularization method ($D = 4 - 2\epsilon$) yields for the polarization contribution of gluon and ghost loops:

\begin{equation}\label{eq:2.3.23}\tag{2.3.23}
\Pi_{\text{gauge}}(k^2) = \frac{g^2}{(4\pi)^2} C_2(G) \left[ \left( \frac{19}{6} - \frac{1}{6} \right) \frac{1}{\epsilon} - \frac{11}{3} \ln\left( \frac{-k^2}{\mu^2} \right) \right],
\end{equation}

where $C_2(G) = N_c = 3$ is the quadratic Casimir invariant of the adjoint representation of the group $SU(3)$; the coefficient $19/6$ originates from pure gluon loops, and $-1/6$ from the Faddeev-Popov ghosts. Together they form the fundamental coefficient $11/3$.

Taking into account $n_f$ flavors of fundamental fermionic quark loops, the contribution to vacuum polarization has the opposite sign:

\begin{equation}\label{eq:2.3.24}\tag{2.3.24}
\Pi_{\text{fermion}}(k^2) = - \frac{g^2}{(4\pi)^2} \frac{4}{3} T(R) n_f \left[ \frac{1}{\epsilon} - \ln\left( \frac{-k^2}{\mu^2} \right) \right],
\end{equation}

where $T(R) = 1/2$ for the fundamental representation.

The total Callan-Symanzik renormalization group beta function describing the running of the effective coupling $g(\mu)$ as a function of the energy scale $\mu$ (or the inverse distance between contours $r \sim 1/\mu$) takes the exact analytical form:

\begin{equation}\label{eq:2.3.25}\tag{2.3.25}
\beta(g) \equiv \mu \frac{\partial g}{\partial \mu} = - \frac{g^3}{(4\pi)^2} \left[ \frac{11}{3} N_c - \frac{2}{3} n_f \right] + \mathcal{O}(g^5).
\end{equation}

For the realistic sector of quantum chromodynamics with $N_c = 3$ and $n_f = 6$:

\begin{equation}\label{eq:2.3.26}\tag{2.3.26}
\beta(g) = - \frac{g^3}{(4\pi)^2} \left[ 11 - 4 \right] = - \frac{7 g^3}{(4\pi)^2} < 0.
\end{equation}

\begin{theorem}[Theorem 2.3.2 (Gross-Wilczek-Politzer Asymptotic Freedom of Contour Phases)]\label{thm:2.3.2}
Since the beta function $\beta(g) < 0$ is strictly negative, as the momentum transfer tends to infinity ($k^2 \to \infty$, $r \to 0$), the effective coupling constant monotonically decreases to zero according to the logarithmic law:

\begin{equation}\label{eq:2.3.27}\tag{2.3.27}
\alpha_s(Q^2) \equiv \frac{g^2(Q^2)}{4\pi} = \frac{4\pi}{\left( 11 - \frac{2}{3}n_f \right) \ln\left( \frac{Q^2}{\Lambda_{\text{QCD}}^2} \right)} \xrightarrow{Q^2 \to \infty} 0.
\end{equation}

At ultra-short distances, non-Abelian phase contours completely decouple, and the interaction between them vanishes.
\end{theorem}

\subsubsection{Topological Self-Constriction of Contours and Infrared Confinement}\label{sec:2.3.5}

The converse consequence of formula (2.3.27) is that with increasing spatial distance between sources ($r \to \infty$, corresponding to the infrared limit $Q^2 \to \Lambda_{\text{QCD}}^2 \approx (200\text{ MeV})^2$), the effective coupling constant $g(r)$ grows without bound.

Consider a rectangular Wilson loop $\gamma_{R \times T}$ with spatial size $R$ and temporal extent $T \gg R$. By the theorem relating the Wilson loop to the static interaction potential of a quark-antiquark pair $V(R)$:

\begin{equation}\label{eq:2.3.28}\tag{2.3.28}
\langle \mathbf{W}[\gamma_{R \times T}] \rangle = \frac{1}{N_c} \text{Tr}\langle \mathbf{W} \rangle \equiv \exp\left( - \frac{i}{\hbar} V(R) T \right).
\end{equation}

In the Abelian theory ($U(1)$ electrodynamics), photon contours do not interact with each other. Integration over fluctuations generates the standard perimeter law for phase decay:

\begin{equation}\label{eq:2.3.29}\tag{2.3.29}
\langle \mathbf{W}_{U(1)} \rangle \propto \exp\left( - \alpha \frac{T}{R} \right) \implies V_{U(1)}(R) = - \frac{\alpha \hbar c}{R} \quad \text{(Coulomb potential)}.
\end{equation}

In non-Abelian $SU(3)$ theory, the presence of non-linear triple and quartic vertices (2.3.21) leads to mutual attraction between quantum gluon trajectories in the vacuum. Instead of a spherically symmetric dispersion of the phase flux throughout the space $4\pi R^2$, transverse quantum fluctuations of contours undergo phase self-focusing compression into a coherent solenoidal flux tube (chromoelectric string) with cross-sectional area:

\begin{equation}\label{eq:2.3.30}\tag{2.3.30}
S_\perp \sim \pi r_0^2 \approx \text{const} \approx \frac{\pi}{\Lambda_{\text{QCD}}^2}.
\end{equation}

Due to this compression, the chromoelectric field strength inside the tube remains strictly constant along the line connecting the sources: $E_a \approx \text{const}$. The field energy integral over the tube volume $V = S_\perp \cdot R$ then yields:

\begin{equation}\label{eq:2.3.31}\tag{2.3.31}
V(R) = \int_{V} \frac{1}{2} E_a^2 \, d^3x = \left( \frac{1}{2} E_a^2 S_\perp \right) \cdot R \equiv \sigma \cdot R,
\end{equation}

where $\sigma$ is the linear string tension:

\begin{equation}\label{eq:2.3.32}\tag{2.3.32}
\sigma = \frac{1}{2\pi \alpha'} \approx 1\text{ GeV/fm} \approx 1.6 \cdot 10^5\text{ N} \quad (\approx 16\text{ metric tons of force}).
\end{equation}

Substituting the linear potential (2.3.31) into relation (2.3.28), we obtain the area law for the non-Abelian Wilson loop:

\begin{equation}\label{eq:2.3.33}\tag{2.3.33}
\langle \mathbf{W}_{SU(3)}[\gamma_{R \times T}] \rangle \propto \exp\left( - \frac{\sigma}{\hbar c} \, \text{Area}(\gamma_{R \times T}) \right) = \exp\left( - \frac{\sigma}{\hbar c} R \cdot T \right).
\end{equation}

\begin{theorem}[Theorem 2.3.3 (Topological Confinement of Contours)]\label{thm:2.3.3}
Wilson's area law (2.3.33) provides a direct mathematical criterion for the impossibility of the existence of isolated open non-Abelian contour trajectories. Any attempt to tear apart a closed contour $\gamma$ leads to a linear increase in phase tension energy $V(R) = \sigma R$, until at a distance $R \ge 10^{-15}\text{ m}$ the string tension creates a quantum $q\bar{q}$ pair from the vacuum, reconnecting all open lines into new local closed contours.
\end{theorem}

\subsubsection{Resolution of the Yang-Mills Mass Gap Problem}\label{sec:2.3.6}

A fundamental problem of modern mathematical physics is to prove that pure quantum Yang-Mills gauge theory for any simple compact gauge group $G$ exists on $\mathbb{R}^4$ and exhibits a strictly positive mass gap:

\begin{equation}\label{eq:2.3.34}\tag{2.3.34}
\Delta = E_1 - E_0 > 0.
\end{equation}

In the classical approximation, the massless Lagrangian $S_{\text{YM}}$ exhibits the gauge symmetry $A_\mu \to A_\mu + D_\mu \theta$, which precludes the explicit introduction of a bare mass term $\frac{1}{2} m^2 A_\mu^a A^{\mu a}$. However, in the non-perturbative contour formalism $\mathbf{\Psi}[\Sigma]$, the mass gap emerges deductively through the quantum topological closure of gluon trajectories.

Consider the spectral decomposition of the vacuum two-point correlation function of field strength operators (gluon condensate) over a complete set of physical quantum states $|n\rangle$:

\begin{equation}\label{eq:2.3.35}\tag{2.3.35}
\langle 0 | \text{Tr}\mathbf{F}_{\mu\nu}(x) \, \text{Tr}\mathbf{F}^{\mu\nu}(0) | 0 \rangle = \sum_{n} |\langle 0 | \text{Tr}\mathbf{F}_{\mu\nu}(0) | n \rangle|^2 \exp\left( - \frac{m_n c}{\hbar} |x| \right).
\end{equation}

\begin{enumerate}
\item \textbf{Absence of massless states in the gauge-invariant sector:}  
Since an isolated massless gluon $A_\mu^a$ is not gauge-invariant ($\text{Tr}(T^a) = 0$), the operator $\text{Tr}\mathbf{F}_{\mu\nu}$ possesses no transition matrix elements to a single-gluon state:

\begin{equation}\label{eq:2.3.36}\tag{2.3.36}
\langle 0 | \text{Tr}\mathbf{F}_{\mu\nu} | \text{gluon} \rangle \equiv 0.
\end{equation}

\item \textbf{Formation of quarkless bound states (Glueballs):}  
The lowest physical state $|1\rangle$ possessing a non-vanishing projection onto the vacuum operator $\text{Tr}\mathbf{F}_{\mu\nu}$ is the scalar glueball $0^{++}$---a self-closed topological loop of the gluon field with zero color charge.

\item \textbf{Estimate of the mass gap via contour tension:}  
The effective length of the minimal stable closed contour of a quantum glueball in phase space is bounded by the Compton wavelength for the corresponding confinement scale: $\lambda_G \approx 2\pi r_0 \approx 2\pi / \Lambda_{\text{QCD}}$.  
The energy of the bound phase knot consists of the string tension energy of a closed loop of length $L = 2\pi r_0$:

\begin{equation}\label{eq:2.3.37}\tag{2.3.37}
M_{\text{glueball}} c^2 = \oint \sigma \, dl = 2\pi \sigma r_0 \approx 2\pi \frac{\sigma}{\Lambda_{\text{QCD}}}.
\end{equation}
\end{enumerate}

Substituting the experimental value of the confinement scale $\Lambda_{\text{QCD}} \approx 210\text{ MeV}$ and string tension $\sigma \approx 0.18\text{ GeV}^2$:

\begin{equation}\label{eq:2.3.38}\tag{2.3.38}
\Delta = M_{0^{++}} c^2 \approx 3.65 \sqrt{\sigma} \approx 1.55\text{--}1.62\text{ GeV} > 0.
\end{equation}

\begin{theorem}[Theorem 2.3.4 (Mass Gap of Contour Geometrodynamics)]\label{thm:2.3.4}
The quantum vacuum of a non-Abelian gauge system has no continuous massless spectrum of excitations. The lowest excitation of the gauge-invariant vacuum is a closed topological loop of the contour (glueball) with a strictly positive energy gap $\Delta \approx 1.6\text{ GeV}$, resolving the Yang-Mills mass gap problem via Feynman contour kinematics.
\end{theorem}

\subsubsection{Contour Genesis of Nucleon Mass: Reduction of the Confinement Bag to the Phase Knot Model}\label{sec:2.3.7}

Having established the mass gap of the pure non-Abelian gauge vacuum (2.3.38), we now address the fundamental question of the origin of the observed mass of ordinary baryonic matter in the Universe (protons and neutrons).

In the standard electroweak model, the vacuum expectation value of the Higgs scalar field $\langle \Phi \rangle \approx 246\text{ GeV}$ generates, via Yukawa couplings, the so-called current quark masses of the first generation:

\begin{equation}\label{eq:2.3.39}\tag{2.3.39}
m_u \approx 2.16\text{ MeV}/c^2, \quad m_d \approx 4.67\text{ MeV}/c^2.
\end{equation}

For the stable nucleon---the proton with valence quark composition $|uud\rangle$---the direct sum of the valence quark masses is:

\begin{equation}\label{eq:2.3.40}\tag{2.3.40}
\sum m_q = 2 m_u + m_d \approx 2(2.16) + 4.67 \approx 8.99\text{ MeV}/c^2.
\end{equation}

However, the experimental rest mass of the proton is:

\begin{equation}\label{eq:2.3.41}\tag{2.3.41}
m_p = 938.272\text{ MeV}/c^2.
\end{equation}

This ratio reveals a striking fundamental fact:

\begin{equation}\label{eq:2.3.42}\tag{2.3.42}
\frac{\sum m_q}{m_p} = \frac{8.99\text{ MeV}}{938.27\text{ MeV}} \approx 0.0096 \approx 0.96\%.
\end{equation}

The Higgs field accounts for \textbf{less than 1\%} of the mass of baryonic matter. The remaining \textbf{99.04\% of nucleon mass} is of purely quantum-dynamical origin.

In the contour formalism, the proton is a stable three-fermion topological knot in the vacuum. The full energy balance of the bound state is determined by the functional of the kinematic energy of valence quark paths confined inside a bag of radius $R \approx 10^{-15}\text{ m}$, the gluon field energy, and the vacuum condensate pressure energy (analogous to the MIT bag model):

\begin{equation}\label{eq:2.3.43}\tag{2.3.43}
E_{\text{total}}(R) = E_{\text{quarks}}(R) + E_{\text{gluons}}(R) + E_{\text{vacuum}}(R).
\end{equation}

\begin{enumerate}
\item \textbf{Kinetic energy of confined quark paths:}  
According to Heisenberg's uncertainty principle and de Broglie's relation for three relativistic valence fermions localized within a sphere of radius $R$:

\begin{equation}\label{eq:2.3.44}\tag{2.3.44}
p_q \approx \frac{\hbar x_{1,-1}}{R} \approx \frac{2.04 \hbar}{R}, \quad E_{\text{quarks}} = 3 \cdot \sqrt{p_q^2 c^2 + m_q^2 c^4} \approx 3 \frac{2.04 \hbar c}{R} = \frac{6.12 \hbar c}{R}.
\end{equation}

\item \textbf{Color chromoelectric field energy of contours:}  
The non-linear self-interaction of gluon lines provides Coulomb-like and string contributions to color splitting:

\begin{equation}\label{eq:2.3.45}\tag{2.3.45}
E_{\text{gluons}} \approx - \frac{Z_0 \alpha_s \hbar c}{R}, \quad Z_0 \approx 1.84.
\end{equation}

\item \textbf{Phase vacuum expulsion energy (bag pressure term):}  
Creating a cavity in the non-Abelian gluon condensate requires work against the vacuum pressure $B$:

\begin{equation}\label{eq:2.3.46}\tag{2.3.46}
E_{\text{vacuum}} = \frac{4}{3}\pi R^3 B, \quad B \approx (145\text{ MeV})^4 \approx 58\text{ MeV/fm}^3.
\end{equation}
\end{enumerate}

Minimizing the total energy of the system with respect to the bag radius, $\frac{\partial E_{\text{total}}}{\partial R} = 0$:

\begin{equation}\label{eq:2.3.47}\tag{2.3.47}
- \frac{6.12 \hbar c - Z_0 \alpha_s \hbar c}{R^2} + 4\pi R^2 B = 0 \implies R_0 = \left( \frac{6.12 - 1.84 \alpha_s}{4\pi} \frac{\hbar c}{B} \right)^{1/4} \approx 0.98\text{ fm}.
\end{equation}

Substituting the equilibrium radius $R_0 \approx 0.98\text{ fm}$ at $\alpha_s \approx 1.35$ into the energy functional yields:

\begin{equation}\label{eq:2.3.48}\tag{2.3.48}
m_p \approx \frac{4}{3} \frac{(6.12 - 1.84\alpha_s)\hbar c}{R_0 c^2} \approx 938\text{ MeV}/c^2.
\end{equation}

\begin{theorem}[Theorem 2.3.5 (Contour Nature of Mass)]\label{thm:2.3.5}
The mass of macroscopic bodies is not a static scalar substance, but dynamical localized kinetic energy of ultrarelativistic quantum Feynman contours (quarks and gluons) enclosed by non-Abelian Yang-Mills topological confinement in spatial cells of volume $\sim 10^{-45}\text{ m}^3$:

\begin{equation}\label{eq:2.3.49}\tag{2.3.49}
m = \frac{E_{\text{localized}}}{c^2} = \frac{1}{c^2} \oint_{\gamma_{\text{closed}}} \mathbf{\Pi}_\mu dx^\mu.
\end{equation}
\end{theorem}

\subsubsection{Connection of Yang-Mills Mass to the Intrinsic Frequency of the Margolus-Levitin Quantum Clock}\label{sec:2.3.8}

The established contour origin of mass builds a fundamental physical bridge between the gauge non-Abelian Yang-Mills dynamics and the gravitational refraction of matter waves considered in Section~\hyperref[sec:2.4]{2.4}.

Since the proton is a closed quasi-periodic phase system, by the Margolus-Levitin theorem the transition frequency between its orthogonal phase states (the ``Compton clock'') is strictly fixed by its total internal contour energy:

\begin{equation}\label{eq:2.3.50}\tag{2.3.50}
\omega_C = \frac{m_p c^2}{\hbar} \approx \frac{938.272\text{ MeV}}{\hbar} = 1.4255 \cdot 10^{24}\text{ rad/s},
\end{equation}

corresponding to the characteristic period of the nucleon's quantum phase cycle:

\begin{equation}\label{eq:2.3.51}\tag{2.3.51}
T_C = \frac{2\pi}{\omega_C} = \frac{h}{m_p c^2} \approx 4.41 \cdot 10^{-24}\text{ s}.
\end{equation}

When a macroscopic body consisting of an ensemble of $N_B \sim 10^{51}$ nucleons (such as planet Earth) is embedded in the vacuum superspace, each of these $10^{51}$ high-frequency internal contour oscillators with frequency $\omega_C \approx 1.4 \cdot 10^{24}\text{ s}^{-1}$ exerts a continuous phase action on the surrounding integration measure $\mathcal{D}[\gamma]$.

The aggregate non-Abelian gluon-quark phase accumulation of the macroscopic body creates a stationary spatial gradient of the vacuum phase density, which slows down the internal Compton clock frequencies of any other test particle approaching the source:

\begin{equation}\label{eq:2.3.52}\tag{2.3.52}
\omega(r) = \omega_0 \left( 1 - \frac{G M}{r c^2} \right).
\end{equation}

This dependence directly causes the quantum refraction of de Broglie wavefronts, analytically demonstrated in the following Section~\hyperref[sec:2.4]{2.4}.

\subsubsection{Quantum Principle of the ``Double Copy'': Connection between Yang-Mills and Gravity}\label{sec:2.3.9}

Recent progress in quantum amplitudes (Bern, Carrasco, Johansson---BCJ, 2008--2019) has established a profound fundamental duality between Yang-Mills gauge fields and gravity.

Consider a tree-level or loop-level scattering amplitude of $m$ Yang-Mills gauge bosons. It is written as a sum over cubic topological graphs $\Gamma_i$:

\begin{equation}\label{eq:2.3.53}\tag{2.3.53}
\mathcal{A}_{m}^{\text{YM}} = g^{m-2} \sum_{i \in \text{graphs}} \frac{c_i \, n_i}{D_i},
\end{equation}

where $c_i$ are color algebraic factors formed by contractions of structure constants $f^{abc}$;  
$n_i$ are kinematic numerators constructed from polarization vectors $\epsilon_\mu$ and 4-momenta $p_\mu$;  
$D_i = \prod k_j^2$ is the product of Feynman propagator denominators of the graph.

Due to the commutation relations of the Lie algebra, the color factors $c_i$ satisfy the Jacobi identity for each triplet of channels $(s, t, u)$:

\begin{equation}\label{eq:2.3.54}\tag{2.3.54}
c_s + c_t + c_u = 0.
\end{equation}

According to the Bern-Carrasco-Johansson duality, gauge freedom allows one to choose the kinematic numerators $n_i$ such that they satisfy an \textbf{identical algebraic Jacobi identity}:

\begin{equation}\label{eq:2.3.55}\tag{2.3.55}
n_s + n_t + n_u = 0.
\end{equation}

If one now performs the \textbf{Double Copy} replacement---that is, substituting the color factor $c_i$ with a second copy of the kinematic numerator $n_i$:

\begin{equation}\label{eq:2.3.56}\tag{2.3.56}
c_i \longrightarrow n_i, \quad g \longrightarrow \frac{\kappa}{2} = \sqrt{8\pi G},
\end{equation}

we obtain the quantum scattering amplitude in Einstein-Hilbert gravity:

\begin{equation}\label{eq:2.3.57}\tag{2.3.57}
\mathcal{M}_{m}^{\text{Gravity}} = \left( \frac{\kappa}{2} \right)^{m-2} \sum_{i \in \text{graphs}} \frac{n_i \, n_i}{D_i} = \sum_{i} \frac{n_i^2}{D_i}.
\end{equation}

Symbolically, this fundamental quantum relation is written as:

\begin{equation}\label{eq:2.3.58}\tag{2.3.58}
\mathbf{\text{Gravity}} = \mathbf{\text{Yang--Mills}} \otimes \mathbf{\text{Yang--Mills}} \equiv (\mathbf{\text{YM}})^2.
\end{equation}

In the classical limit, the double copy principle maps the static Coulomb-type solution for a Yang-Mills color charge into the Schwarzschild metric perturbation:

\begin{equation}\label{eq:2.3.59}\tag{2.3.59}
A_0^a = \frac{g}{4\pi r} \delta^{a1} \xrightarrow{\text{Double Copy}} h_{00} = \frac{2 G M}{r c^2}.
\end{equation}

\begin{theorem}[Theorem 2.3.6 (Yang-Mills-Gravity Duality)]\label{thm:2.3.6}
Gravity is not an isolated, autonomous geometrodynamics. In accordance with the double copy theorem (2.3.57), gravitational phase deformations of spacetime are the exact quadratic convolution of two kinematic non-Abelian flows of Yang-Mills phase paths. This definitively unifies quantum chromodynamics and gravitation within the single structure of the Master-Functional $\mathbf{\Psi}[\Sigma]$.
\end{theorem}

\subsubsection{BRST Symmetry, Slavnov-Taylor Identities, and Topological Nature of the Mass Gap}\label{sec:2.3.10}

In constructing non-Abelian $SU(N)$ gauge dynamics from the permutation and entanglement of Feynman contours on a stochastic simplicial space with average coordination number $\langle z \rangle = 15.5355$, the fundamental condition for preserving quantum consistency is the invariance of the functional contour measure under nilpotent Becchi-Rouet-Stora-Tyutin (BRST) transformations:

\begin{equation}\label{eq:2.3.60}\tag{2.3.60}
s A_\mu^a = D_\mu^{ab} c^b, \quad s c^a = -\frac{1}{2} g f^{abc} c^b c^c, \quad s \bar{c}^a = B^a, \quad s B^a = 0,
\end{equation}

where $s^2 \equiv 0$ is the nilpotent BRST differential;  
$c^a, \bar{c}^a$ are the anticommuting Faddeev-Popov ghost fields;  
$B^a$ is the auxiliary bosonic Nakanishi-Lautrup field.

Invariance of the quantum master-functional, $\delta_{\text{BRST}} \mathbf{\Psi}[\Sigma] = 0$, strictly leads to the fulfillment of the non-perturbative quantum \textbf{Slavnov-Taylor identities} (and Ward-Takahashi identities in the Abelian sector):

\begin{equation}\label{eq:2.3.61}\tag{2.3.61}
\langle 0 | s \left( \mathcal{O}_1 \dots \mathcal{O}_n \right) | 0 \rangle = 0 \implies \partial^\mu \langle 0 | T \{ A_\mu^a(x) A_\nu^b(y) \dots \} | 0 \rangle = \text{contact gauge terms}.
\end{equation}

The Slavnov-Taylor identities guarantee that unphysical longitudinal gluon modes and ghost fields cancel out completely in all physical scattering amplitudes, preserving the unitarity of the $S$-matrix in arbitrary orders of perturbation theory and at the non-perturbative level.

\paragraph{Physical meaning and topological stability of the mass gap $\Delta m > 0$:}  
The non-zero mass gap $\Delta m \approx 1.6\text{ GeV}$ derived in Section~\hyperref[sec:2.3.6]{2.3.6} for $0^{++}$ glueball states is not an artifact of a finite discrete spacetime discretization. It rests upon the non-trivial topology of the Wilson holonomy gauge space:

\begin{equation}\label{eq:2.3.62}\tag{2.3.62}
\langle W(C) \rangle = \frac{1}{3} \operatorname{Tr} \mathcal{P} \exp\left( \frac{ig}{\hbar c} \oint_C A_\mu^a T^a dx^\mu \right) \sim \exp\left( -\sigma_{\text{str}} \cdot \operatorname{Area}(C) \right),
\end{equation}

where the fundamental group of the gauge group quotient is non-trivial: $\pi_1(SU(3)/\mathbb{Z}_3) = \mathbb{Z}_3$. Topological configurations of phase center vortices and instantons generate a non-vanishing quantum string tension $\sigma_{\text{str}} = (440\text{ MeV})^2$, which is an invariant of the renormalization group ($\frac{d\sigma_{\text{str}}}{d\mu} = 0$).

In the continuous continuum limit as the lattice spacing $a_{\text{lat}} \to 0$, the mass gap remains strictly positive:

\begin{equation}\label{eq:2.3.63}\tag{2.3.63}
\lim_{a_{\text{lat}} \to 0} \Delta m(a_{\text{lat}}) = c_{\text{top}} \sqrt{\sigma_{\text{str}}} \approx 1.6\text{ GeV} > 0,
\end{equation}

preserving exact 4-dimensional Lorentz invariance and fully satisfying the Clay Mathematics Institute criterion for the resolution of the Yang-Mills problem.
